\chapter{An evidence program for the thesis}\label{ch:evidence}

The empirical work of Chapters~\ref{ch:p1} and~\ref{ch:p2} tests each paper's own predictions. This chapter states what that evidence establishes for the thesis as a whole, what it does not, and two tests designed to test the thesis directly. Both tests are specified here and have not been run.

\section{What the existing evidence establishes}\label{sec:ev_existing}

\paragraph{Prices do not reveal the agents.} In 30 Binance perpetuals, trader-class positions pin down memory rates 28 to 124 times more precisely than prices and flow, and in a pre-registered test the CFTC's trader-class reports do the same in all eight regulated futures markets examined (Sections~\ref{p1:sec:disc_data} and~\ref{p1:sec:cot}). This is the identification half of the thesis in data: the merged measure leaves class memories essentially unrestricted.

\paragraph{The safety statistic carries information.} In thin Binance books, six-hour estimates of the amplification margin forecast the size of the next hour's adverse excursion out of sample in every pre-registered forecasting test, beyond a full volatility-forecasting baseline and VPIN; the effect is about a fifth as large in thin US stocks, negligible on Bybit and absent on Coinbase (Section~\ref{p2:sec:thin} and Table~\ref{p2:tab:alltests}).

\paragraph{What it does not establish.} The forecasting result is reduced form. Tests of a leverage or liquidation mechanism fail (Section~\ref{p2:sec:p22}), and the follow-up that points to slow reversion is close to a check of how the statistic is built (Section~\ref{p2:sec:p25}). As a risk control, the certified guard improves on a static limit but not, at matched exposure, on a volatility cap (Section~\ref{p2:sec:p24}). None of the existing tests checks the thesis's central claim directly: that the certificate computed without identification gets safety right while identification fails.

\section{Test E1: the certificate against ground truth}\label{sec:ev_e1}

\paragraph{Idea.} Where agents' classes are disclosed, the total sensitivity can be computed twice: once from prices and aggregate flow alone, as a certificate, and once by summing class-level estimates, $\hat\Phi_{\rm class}=\sum_c\hat w_c/(1-\hat r_c)$, as the disclosed benchmark. The thesis predicts that the two agree within their intervals, while the class memories themselves are not recoverable from prices, which Chapter~\ref{ch:p1} has already shown for the same markets.

\paragraph{Design.} Markets: the eight CFTC markets of Section~\ref{p1:sec:cot} (weekly) and the 30 Binance perpetuals of Section~\ref{p1:sec:disc_data} (five-minute), on dates not used before. Certificate: the delta-method interval for $\Phi$ from the price-and-flow response kernel with a declared memory bound, as in Chapter~\ref{ch:p2}. Benchmark: the class-composed $\hat\Phi_{\rm class}$ with a block-bootstrap interval. Primary 1: the two intervals overlap in at least 75\% of markets. Primary 2: in the same markets, the price-only intervals for the individual class memories are at least ten times wider than the class-based ones. Caveat registered in advance: disclosed classes need not exhaust the agents, so disagreement can mean either that the certificate is wrong or that unclassified agents matter; the test reports which classes account for the gap.

\pendbox{E1}{Specification to be written, hashed and pushed to the public repository on its own before any test data are downloaded; results reported whatever they show.}

\section{Test E2: a prospective test of the mechanism}\label{sec:ev_e2}

\paragraph{Idea.} Every pre-registered test so far used archived data. A prospective test removes any doubt about the order of registration and data: the data do not exist when the specification is written. Binance publishes forced liquidations on a public market-data stream, at most one per symbol each second, so the stream is a sample of liquidations rather than the full record, and its history is not archived publicly. Recording the stream after registration gives liquidation data no one could have seen; the sampling is the same in every hour and is stated as a limitation in advance.

\paragraph{Design.} Record the forced-liquidation stream and order-book depth for the thin perpetuals of Chapter~\ref{ch:p2} for at least three months after the specification is pushed. Primary 1: six-hour $\hat C_3$ forecasts next-hour sampled liquidation volume beyond the volatility, impact and lagged-liquidation baseline. Primary 2: adding realized next-hour liquidation volume to the excursion model removes at least half of $\hat C_3$'s forecasting gain. Either outcome is informative: a pass identifies the channel that the open-interest proxy of Part~22 could not; a fail, with direct liquidation data, closes it.

\pendbox{E2}{Recording infrastructure, specification and power calculation to be completed and the specification pushed before recording starts.}

\section{Rules for every new test}\label{sec:ev_rules}

\begin{enumerate}
\item The specification, its hash and an unlock counter are pushed to the public repository in a commit of their own, so the platform records when they became public, before any test data are downloaded or recorded.
\item Each test has at most two primaries with a pre-stated pass rule and a pre-stated correction across them.
\item A failed primary is reported as failed. An explanation found afterwards is labelled post hoc, and any test that follows from it is labelled as registered after the failure.
\item Every number in the thesis is produced by code that writes a result file, and an automated check compares the printed numbers with those files.
\end{enumerate}
